\section{Methodology}
\label{sec:method}
The motivation (\S\ref{sec:motivation}) established two facts: understanding
saturates at a mid-stack depth and is therefore cacheable in a single hidden
state, while the top layers perform the query-conditioned generation and must
be recomputed with the query present. CoMem turns these facts into a system by
partitioning the transformer along the \emph{depth} axis rather than the
\emph{token} axis: it caches one mid-layer hidden per chunk at write time, and
at read time retrieves a fixed number of these hiddens and recomputes only the
upper layers. Read compute and memory are then constant in context length. We
first fix notation, then describe the Write and Read passes, the retrieval
(selector) module, the split-depth knob $j$, and the optional self-distillation
that deepens the cacheable point.

\paragraph{Formal notation.} Let a decoder-only backbone have $L$ transformer
layers, hidden width $d$, and layers indexed $[0{:}L]$. CoMem fixes a single
\emph{split depth} $j\in\{0,\dots,L\}$ and a \emph{chunk size} $c{=}512$. For a
chunk of $c$ tokens, $h_j\in\mathbb{R}^{c\times d}$ denotes its depth-$j$
residual-stream hidden state---the sole quantity cached per chunk. Because a
full KV cache stores a key and a value tensor at \emph{every} one of the $L$
layers, whereas CoMem stores one hidden at a \emph{single} depth, its per-token
storage is $\approx 1/(2L)$ of a full-depth KV cache (on Qwen3-8B, $L{=}36$, so
$\approx 1/72$). At read time we assemble a \emph{pack}
$\texttt{[sink}\,;\,\texttt{selected }h_j\,;\,\texttt{query }h_j\texttt{]}$
built from the top-$k$ retrieved chunk hiddens; its length is
$\text{sink}+k\!\cdot\!c+\text{query}\approx 6.7\text{k}$ tokens for $k{=}12$,
$c{=}512$, and---crucially---is \emph{independent of the context length}.

\paragraph{Write.} The document is streamed and split into disjoint
$c$-token chunks. Each chunk is encoded through the bottom layers $[0{:}j]$
under \emph{chunk-local} RoPE positions (each chunk starts at position $0$), and
its depth-$j$ residual-stream hidden $h_j$ is written to an external store keyed
by the chunk's raw text. Only the bottom $j$ layers run and only one hidden per
token is retained, so the write pass is shallow ($O(L_{\text{ctx}})$ in the
number of tokens but a factor $L/j$ cheaper than a full forward) and the store
grows at $\approx 1/(2L)$ the rate of a KV cache. Chunk-local positions make
every $h_j$ position-canonical, so cached hiddens are reusable at read time
regardless of where the chunk originally sat in the document.

\paragraph{Read.} Given a query, the selector (below) returns the indices of the
top-$k$ chunks. We gather their stored hiddens, prepend a small attention
\emph{sink} and append the query's own $h_j$ (obtained by running the query text
through $[0{:}j]$), and assign the whole pack \emph{fresh}, contiguous RoPE
positions. We then recompute layers $[j{:}L]$ over this pack with
\emph{full cross-pack attention}---every token attends to every other token in
the pack---to obtain next-token logits. Full attention is deliberate and
load-bearing: reusing each chunk's cached upper-layer attention block-diagonally
(as in KV-reuse pipelines such as CacheBlend/RAGCache \citep{cacheblend,ragcache})
would prevent the query and the retrieved facts from interacting and collapses
multi-fact disambiguation; we motivate this in \S\ref{sec:motivation} and ablate
it in \S\ref{sec:experiments} (Table~\ref{tab:crosschunk}). Because the pack
length is fixed by $k$ and $c$, both the read FLOPs and the read-time KV memory
are constant in context length.

\subsection{Chunk selection}
\label{subsec:selector}
Retrieval is external, non-differentiable, and---by default---requires no model
forward pass. This keeps the memory a plain, non-evolving store of hiddens and
makes the selector a swappable component.

\paragraph{Default: one-shot BM25.} The default selector scores all context
chunks against the query with BM25 lexical matching over the shared chunk
tokenization and returns the top-$k$. It is a single pass, purely CPU-side, and
adds negligible latency. For entity-needle workloads BM25 is near-optimal:
given the same read pack, it reaches within a few points of an oracle that
always includes the gold chunk (\S\ref{sec:experiments},
Table~\ref{tab:selector}), so the practical long-range accuracy of CoMem on
these tasks is bounded by retrieval recall, not by any information lost in the
cached hidden.

\paragraph{Iterative BM25 for reference chains.} Single-shot lexical retrieval
fails on multi-hop reference chains (e.g.\ variable tracking, where
$\texttt{VAR3}\!\to\!\texttt{VAR1}\!\to\!\texttt{VAR2}{=}42$): the query
matches the last link but not the chunks that define the intermediate
variables. We introduce \texttt{iter\_bm25}, a forward-free chain-following
variant. Round~1 is identical to one-shot BM25. Each subsequent round
re-queries BM25 using the concatenated token text of the previous round's
selected chunks, so a freshly retrieved chunk's high-IDF variable name pulls in
the chunk that \emph{defines} that variable---one hop backward per round. We
retrieve \texttt{iter\_hop\_topk} chunks per round for
$\lceil k/\texttt{iter\_hop\_topk}\rceil$ rounds, skipping candidates with zero
lexical overlap as a chain-end guard. \texttt{iter\_bm25} adds \emph{no}
training and \emph{no} extra model forward---only additional CPU-side BM25
lookups over the same corpus---yet it recovers variable tracking from the
$20$--$28$ range to $92$--$97$ and, unlike single-pass selectors, does not
degrade with context length (\S\ref{sec:experiments},
Table~\ref{tab:itervt}). It is scoped to reference-chain tasks: on
independent-needle tasks it is a no-op (single-needle $100{=}100$) and can
slightly dilute multi-key retrieval, so the default single-shot BM25 remains the
selector elsewhere.

\paragraph{Comparison selectors.} Two further selectors and an upper bound are
used only for analysis (\S\ref{sec:experiments}, Table~\ref{tab:selector}):
\texttt{reader\_attn}, which scores chunks by the cosine similarity between the
query's $h_j$ and each chunk's $h_j$; \texttt{recency}, which keeps the most
recent chunks (a StreamingLLM-style positional prior); and \texttt{oracle},
which always includes the gold chunk and quantifies the retrieval ceiling. None
is used in the default CoMem configuration.

\subsection{The depth knob and training}
\label{subsec:jknob}

\paragraph{The $j$ knob.} The split depth $j$ interpolates between two exact
endpoints. At $j{=}0$ the read recomputes \emph{all} layers over the retrieved
raw tokens; a self-test confirms this equals the full forward to
$\max|\Delta\text{logit}|<10^{-4}$ on the dense 8B backbone (and exactly
$0.0$ on the Hunyuan MoE, \S\ref{sec:experiments}), so $j{=}0$ is exact
full-context RAG recompute---the training-free residual-stream recompute of
KV-Direct \citep{kvdirect}. At $j{=}L$ nothing is recomputed and the model is
closed-book. Intermediate $j$ trades cache compressibility (deeper $j$ stores a
more digested, cheaper-to-recompute hidden) against readout fidelity (deeper $j$
has committed more to a query-blind continuation). The cacheable sweet spot
lands at $\approx 0.33$--$0.40\!\cdot\!L$ across backbones---Qwen3-8B at
$j{=}12$ and the Hunyuan 80-layer MoE at $j{\approx}32$
($0.40\!\cdot\!L$)---and $j$ is exposed as a single deployment-time dial rather
than a fixed architectural choice.

\paragraph{Self-distillation.} With no training, the usable cache depth on
Qwen3-8B saturates around $j{\le}9$; recomputing from $j{=}12$ zero-shot loses
accuracy. We recover it with a light self-distillation: a student that reads
from $j{=}12$ is trained to match a $j{=}0$ teacher (which sees the full
recompute) by KL divergence on their next-token distributions. Training uses a
LoRA adapter ($r{=}32$, $\alpha{=}64$) on the backbone, $4000$ steps of 8-GPU
DDP, over plain PG19 continuous windows only---\emph{no synthetic
long-context data, no task supervision, and no retrieval in the loop}. This
alone pushes the usable cache depth from $j{\le}9$ to $j{\ge}12$, and because
the objective is generic language modeling the resulting memory transfers
zero-shot across benchmarks and backbone families.

\paragraph{No special pretraining is required.} We stress that CoMem operates on
a stock, frozen backbone: the depth partition, the retrieval readout, and the
$j$ knob all work zero-training, and the self-distilled LoRA above is a light,
optional add-on rather than a prerequisite. A separate line of work asks whether
the cache point can be made \emph{compressible by design}---inserting a
low-rank funnel at layer $j$ during from-scratch pretraining so the depth-$j$
hidden is intrinsically lower-rank---which is complementary to, but not required
by, CoMem and orthogonal to train-time token-axis compression such as
\citet{kvcat}; we treat it as companion work and make no such claim here.
